\documentclass[11pt,twoside,twocolumn]{usgsreport}
\usepackage{usgsfonts}
\usepackage{usgsgeo}
\usepackage{usgsreporta}

\usepackage{graphicx}
\usepackage{natbib}
\usepackage{enumitem}
\usepackage[T1]{fontenc}
\usepackage{etoolbox}
\usepackage{fancyvrb}

\usepackage{hyperref}
\hypersetup{
    pdftitle={triangle-utils User Guide},
    pdfauthor={triangle-utils developers},
    pdfsubject={triangle-utils User Guide},
    pdfkeywords={triangle-utils, Triangle, triangular grid, Voronoi grid, MODFLOW 6, MODFLOW-USG},
    bookmarksnumbered=true,
    bookmarksopen=true,
    bookmarksopenlevel=1,
    colorlinks=true,
    allcolors={blue},
    pdfstartview=Fit,
    pdfpagemode=UseOutlines,
    pdfpagelayout=TwoPageRight
}

\graphicspath{{./Figures/}}

% the version and the values the text quotes, written by make_figures.py
\input{values.tex}

\renewcommand{\cooperator}{}
\renewcommand{\reporttitle}{triangle-utils: User Guide}
\renewcommand{\coverphoto}{cover.png}
\renewcommand{\GSphotocredit}{Triangular mesh of an example valley created with triangle-utils.}
\renewcommand{\reportseries}{}
\renewcommand{\reportnumber}{}
\renewcommand{\reportyear}{2026}
\ifdef{\reportversion}{\renewcommand{\reportversion}{Version \tuversion}}{}
\renewcommand{\theauthors}{triangle-utils developers}
\renewcommand{\thetitlepageauthors}{Version \tuversion}
\renewcommand{\reportcitingtheauthors}{triangle-utils developers}
\renewcommand{\colophonmoreinfo}{}
\renewcommand{\reportbodypages}{}
\urlstyle{rm}
\renewcommand{\reportwebsiteroot}{}
\renewcommand{\reportwebsiteremainder}{}
\ifdef{\usgsissn}{\renewcommand{\usgsissn}{}}{}
\renewcommand{\theconventions}{}
\definecolor{coverbar}{RGB}{0, 112, 60}
\renewcommand{\bannercolor}{\color{coverbar}}
\renewcommand{\reportrefname}{References Cited}

\begin{document}

% the class cover without the banner, the agency logo (a blank page in place
% of the logo file), and the department and agency lines
{\renewcommand{\coverruleoddpage}{}\renewcommand{\USGSvisidname}{blank}%
 \renewcommand{\doiagency}{}\renewcommand{\textusgs}{}%
 \ifdef{\makefrontcoveralt}{\makefrontcoveralt}{\makefrontcover}}
\ifdef{\makefrontmatterabv}{\makefrontmatterabv}{\makefrontmatter}

\onecolumn
\pagestyle{body}
\RaggedRight
\setlength{\parindent}{2em}

\section*{Abstract}
\addcontentsline{toc}{section}{Abstract}

triangle-utils is a set of utilities for building two-dimensional triangular meshes for groundwater models, such as MODFLOW~6 \citep{modflow6gwf}, from the boundaries, rivers, lakes, and wells of a model area. The utilities use Triangle \citep{shewchuk1996triangle}, which builds meshes whose triangles have no angle smaller than a specified limit. This guide describes generating a mesh that follows the boundary, rivers, and lakes of a model area and has smaller triangles around pumping wells, and refining a mesh by area and by angle, so that no triangle has an angle over a limit, which the flow calculation between the cells of a triangular grid needs. It also describes improving a mesh by moving its vertices, so that fewer triangles have an angle over 90 degrees, building from a mesh a Voronoi grid, whose cells are centered on the mesh's vertices, for MODFLOW 6, and giving the grid layers, with the input of MODFLOW 6's vertex (DISV) and unstructured (DISU) discretizations. Generating, refining, and improving a mesh, and building and layering its grid, are demonstrated on an example valley with a river, a lake, and three pumping wells.

\section{Introduction}

This guide describes triangle-utils version \tuversion, which is built with Triangle version \triangleversion\ from the MODFLOW-ORG/triangle repository, as of its release \trianglerelease\ or later. The utilities generate a triangular mesh from the features of a model area, refine a mesh by area and by angle, improve it by moving its vertices, build a Voronoi or triangle grid from it, map rivers, boundaries, and wells to the grid's cells, and give the grid layers, with the DISV and DISU input of MODFLOW~6. Later versions will write the discretization files of MODFLOW~6 and MODFLOW-USG from the command line, and this guide will describe each addition.

A groundwater model that uses a triangular or Voronoi grid needs a mesh whose triangles follow the model boundary and the features that control flow, such as rivers and lakes, and are small where heads change rapidly, such as near pumping wells. The shape of the triangles matters as well. The flow between two cells of an unstructured grid is calculated from the head difference between the cell centers, which is accurate when the line between the centers is perpendicular to the face the cells share and crosses it. For a triangular grid whose cell centers are the circumcenters of the triangles, the line between two centers is always perpendicular to the shared face, but the circumcenter of a triangle with an angle over 90 degrees lies outside the triangle, so its cell center is outside its cell and the line need not cross the face.

\section{Triangle and Its Terms of Use}

Triangle generates Delaunay triangulations and quality meshes, in which no angle is smaller than a specified limit, by Delaunay refinement \citep{shewchuk2002delaunay}. triangle-utils builds Triangle from the MODFLOW-ORG/triangle repository, which distributes Triangle with the changes needed to build it on current compilers and to return its errors to the calling program instead of ending it. Triangle is copyrighted by its author and may be distributed only free of charge, with its copyright notice and with its source code available; distributing it as part of a commercial system requires an arrangement with the author. These terms apply to every program and library built from triangle-utils, since each contains Triangle. The code of triangle-utils itself is dedicated to the public domain under Creative Commons CC0 1.0. The command \texttt{triangle-utils --license} prints both.

\section{Generating a Mesh}

A mesh is generated from the features of the model area, each given by the coordinates of its points:

\begin{itemize}
  \item polygons, which bound the model area and the areas inside it, such as a lake, that the mesh follows;
  \item hole polygons, whose inside is not meshed, and points inside holes bounded by polygons;
  \item open lines that the mesh follows, such as rivers;
  \item regions, each given by a point inside it, which give the triangles of the region an attribute and a largest triangle area;
  \item points that the mesh includes as vertices, such as observation wells; and
  \item rings of finer triangles around a point, such as a pumping well, made of a vertex at the point, a polygon of points around it at a given radius, and a region inside the polygon with its own attribute and largest triangle area.
\end{itemize}

Points are added along each polygon and line so that no two successive points are farther apart than a spacing given for that feature, which sets the size of the triangles along it. Points that appear in more than one feature are merged into one vertex, and a point of one feature that lies on an edge of another, such as the end of a river on the model boundary between two of its points, divides that edge, so the two features share the point. Points whose coordinates differ only in their last digits, by less than $10^{-12}$ of the largest coordinate (\valleytolerance~micrometers in the example valley), are merged in the same way, and a larger merge distance can be given, for features digitized separately whose shared points do not match exactly. The mesh covers the polygons, without the area inside a hole polygon or anything in it, and its triangles have no angle smaller than a given minimum angle, up to 34 degrees, and no area larger than a given largest area. Triangle spreads the attribute and area of a region only up to the polygons and lines around it, so each part that rivers and other lines divide the model area into needs a region of its own; the example valley's river and tributary divide its aquifer into three.

Figure~\ref{fig:generated} shows a mesh of an example valley of \valleyarea\ square kilometers (km$^2$), with points every \boundaryspacing\ meters (m) along its boundary, every \riverspacing~m along the river and its tributary, and every \lakespacing~m around the lake. Triangles are no larger than \maximumarea\ square meters (m$^2$) in the aquifer and \lakearea~m$^2$ in the lake, and have no angle smaller than \minimumangle\ degrees. Each of the three wells is at the center of a ring with a radius of \wellradius~m and \wellpoints\ points, in which triangles are no larger than \wellarea~m$^2$. The river starts and ends on the boundary, between two of its points, and the tributary starts on it and ends on the river. Figures~\ref{fig:generated}\textit{B}--\textit{D} show the mesh in square windows \wellwindow~m, \riverwindow~m, and \lakewindow~m on a side around the well where the generated mesh has the most triangles with an angle over \maximumangle\ degrees (fig.~\ref{fig:refined}), at the confluence, where the tributary joins the river, and on the lake's shore. The mesh has \generatedvertices\ vertices and \generatedtriangles\ triangles.

\begin{figure}[htbp]
  \centering
  \includegraphics{generated_mesh.pdf}
  \caption{Mesh generated for the example valley, colored by region. \textit{A}, the valley, with the river and its tributary, the lake, the rings around the three wells, and the windows of panels \textit{B}--\textit{D} (boxes). \textit{B}, the mesh around the well at x = \wellx~m and y = \welly~m. The ring around each well gives it triangles of no more than \wellarea~m$^2$ while the triangles around the ring grow to the size of those in the aquifer. \textit{C}, the mesh at the confluence, whose triangles are smallest along the streams and grow away from them. \textit{D}, the mesh on the lake's shore, whose triangles are no larger than \lakearea~m$^2$ in the lake}
  \label{fig:generated}
\end{figure}

\section{Refining a Mesh}

Refining splits the triangles of an existing mesh, keeping its boundary and its polygons and lines. A mesh is refined so that each triangle is no larger than an area given for it, so that every triangle is no larger than a largest area, or so that no angle is smaller than a minimum angle or larger than a maximum angle. The area of each triangle can be set from a model quantity, such as a hydraulic gradient from an earlier simulation.

Refining by a maximum angle reduces the obtuse triangles of a mesh. Triangle can reduce the largest angle to about 95 to 100 degrees, but not to 90 degrees, at which refining does not finish. Refining by a maximum angle stops after a given number of triangles have been split, and the triangles left with an angle over the limit can be counted. The attributes and area limits of regions are kept exactly only where polygons or lines bound them; elsewhere they are averaged with those of neighboring triangles.

Figure~\ref{fig:refined} shows the mesh of Figure~\ref{fig:generated} refined to a maximum angle of \maximumangle\ degrees. In the generated mesh, \generatedover\ triangles have an angle over \maximumangle\ degrees, the largest \generatedlargest\ degrees. The refined mesh has \refinedtriangles\ triangles, of which \refinedover\ have an angle over \maximumangle\ degrees, the largest \refinedlargest\ degrees.

\begin{figure}[htbp]
  \centering
  \includegraphics{refined_mesh.pdf}
  \caption{Triangles with an angle over \maximumangle\ degrees (red) in the mesh of the example valley, with the windows of figures~\ref{fig:generated}\textit{B}--\textit{D} and~\ref{fig:closeups} (boxes). \textit{A}, the mesh generated with a minimum angle of \minimumangle\ degrees (fig.~\ref{fig:generated}). \textit{B}, the same mesh refined to a maximum angle of \maximumangle\ degrees, which leaves \refinedover\ triangles over it}
  \label{fig:refined}
\end{figure}

Figure~\ref{fig:closeups} shows the windows of Figure~\ref{fig:generated} before and after refining. In the generated mesh, \wellwindowover\ triangles around the well, \riverwindowover\ at the confluence, and \lakewindowover\ on the lake's shore have an angle over \maximumangle\ degrees; in the refined mesh, \wellwindowrefinedover, \riverwindowrefinedover, and \lakewindowrefinedover.

\begin{figure}[htbp]
  \centering
  \includegraphics{close_ups.pdf}
  \caption{Triangles with an angle over \maximumangle\ degrees (red) in the windows of figures~\ref{fig:generated}\textit{B}--\textit{D}. \textit{A}, around the well, generated. \textit{B}, around the well, refined. \textit{C}, at the confluence, generated. \textit{D}, at the confluence, refined. \textit{E}, on the lake's shore, generated. \textit{F}, on the lake's shore, refined}
  \label{fig:closeups}
\end{figure}

\section{Improving a Mesh}

Improving moves the vertices of a mesh, other than those on its boundary, polygons, and lines, those of point features and wells, and those between regions, and redraws the edge two triangles share wherever that makes their smallest angles larger. A vertex is moved toward the average of its neighboring vertices, or toward the centroids of its triangles, only where every triangle around it keeps a positive area within its area limit, no angle falls below the smallest angle of the mesh, and the largest angle around it decreases toward a given maximum angle. With a maximum angle of 90 degrees, improving reduces the obtuse triangles that refining leaves, whose circumcenters lie outside them. An obtuse angle at a point on the boundary where only two triangles meet cannot be removed by moving vertices, because the two angles there sum to 180 degrees; refining that part of the boundary adds triangles that can.

By default, no move may make an angle smaller than the smallest angle of the mesh, so improving never makes that angle smaller; some obtuse triangles can be removed only by making an angle smaller, and a smaller minimum angle can be given to allow it.

Figure~\ref{fig:improved} shows the refined mesh of Figure~\ref{fig:refined} improved toward a maximum angle of \improveangle\ degrees. The refined mesh has \refinedoverninety\ triangles with an angle over \improveangle\ degrees and a smallest angle of \refinedsmallest\ degrees. Improved with that smallest angle kept, the mesh has \improvedoverninety\ triangles over \improveangle\ degrees: \improvedboundaryover\ with the angle at a point on the boundary where two triangles meet, \improvedfixedover\ at another point that does not move (on the river, the lake, or a ring), and \improvedfreeover\ at a vertex that moves. Improved with a minimum angle of \tradedangle\ degrees, the mesh has \tradedoverninety\ triangles over \improveangle\ degrees, \tradedfreeover\ of them at a vertex that moves, and a smallest angle of \tradedsmallest\ degrees, so keeping the smallest angle is what left the \improvedfreeover\ at vertices that move.

\begin{figure}[htbp]
  \centering
  \includegraphics{improved_mesh.pdf}
  \caption{Triangles with an angle over \improveangle\ degrees (red) in the mesh of the example valley. \textit{A}, the mesh refined to a maximum angle of \maximumangle\ degrees (fig.~\ref{fig:refined}\textit{B}). \textit{B}, the same mesh improved toward a maximum angle of \improveangle\ degrees, keeping its smallest angle. \textit{C}, the same mesh improved with a minimum angle of \tradedangle\ degrees. Boxes are the windows of figure~\ref{fig:improvedcloseups}}
  \label{fig:improved}
\end{figure}

Figure~\ref{fig:improvedcloseups} shows the windows of Figure~\ref{fig:generated} before and after improving. In the refined mesh, \wellwindowrefinedoverninety\ triangles around the well, \riverwindowrefinedoverninety\ at the confluence, and \lakewindowrefinedoverninety\ on the lake's shore have an angle over \improveangle\ degrees; in the improved mesh, \wellwindowimprovedover, \riverwindowimprovedover, and \lakewindowimprovedover, although the vertices on the shore do not move.

\begin{figure}[htbp]
  \centering
  \includegraphics{improved_close_ups.pdf}
  \caption{Triangles with an angle over \improveangle\ degrees (red) in the windows of figures~\ref{fig:generated}\textit{B}--\textit{D}. \textit{A}, around the well, refined. \textit{B}, around the well, improved. \textit{C}, at the confluence, refined. \textit{D}, at the confluence, improved. \textit{E}, on the lake's shore, refined. \textit{F}, on the lake's shore, improved}
  \label{fig:improvedcloseups}
\end{figure}

\section{Building a Grid for MODFLOW 6}

MODFLOW 6 calculates the flow between two cells from the difference in head between their centers, which represents the flow across their shared face exactly when the line between the centers is perpendicular to that face. In a Voronoi grid, each cell surrounds one vertex of the mesh and contains the points closer to that vertex than to any other, so the face two cells share is perpendicular to the line between their vertices, and the vertices are the cell centers. The corners of a cell are the circumcenters of the triangles around its vertex, the centers of the circles through each triangle's three corners. A cell on the boundary is closed by the midpoints of the boundary's edges, and by its vertices where it turns, so its center is on its edge, and it is not convex where the boundary turns inward, as at a corner of a hole. Each circumcenter must lie inside the model area, and the circumcenters around a vertex must follow one another around it, which holds when no circle through a triangle's corners contains another vertex and no angle opposite the boundary is over 90 degrees. Triangle's conforming Delaunay switch (D) adds vertices until both hold, so the mesh is generated and refined with it.

Every line and polygon inside the model area runs through cell centers, since its vertices are vertices of the mesh, and the cells on it extend to both sides. Along each stream, a cell's center, where MODFLOW 6 calculates its head, is on the channel. A cell on the edge of a region, such as the lake's shore or a well's ring, takes the region of the side with most of its area; of the \shorecells\ cells on the lake's shore, \shorelakecells\ take the lake's region (fig.~\ref{fig:voronoi}\textit{D}).

Figure~\ref{fig:voronoi} shows the Voronoi grid of the example valley, built from a mesh generated and refined as in Figures~\ref{fig:generated} and~\ref{fig:refined} with conforming Delaunay triangles, \delaunaytriangles\ in all. The grid has \voronoicells\ cells and \voronoivertices\ cell corners.

\begin{figure}[htbp]
  \centering
  \includegraphics{voronoi_grid.pdf}
  \caption{Voronoi grid of the example valley, colored by region, with the windows of figures~\ref{fig:generated}\textit{B}--\textit{D} (boxes). \textit{A}, the valley. \textit{B}, around a well. \textit{C}, at the confluence. \textit{D}, on the lake's shore, where each cell takes the region of most of its area. The cell centers (dots) are the vertices of the mesh}
  \label{fig:voronoi}
\end{figure}

A grid can also be made of the triangles themselves, centered on their centroids or their circumcenters. A circumcenter also gives faces perpendicular to the lines between centers, but it lies outside an obtuse triangle, so the obtuse triangles are refined and improved first. Either grid is written as the vertices and cells of a MODFLOW 6 DISV file.

The unstructured discretization (DISU) of MODFLOW 6 and MODFLOW-USG describes the grid by its connections, with its vertices optional in MODFLOW 6: for each face two cells share, the distance from each cell's center to the face, measured along the face's normal (CL12), the face's width (HWVA), and the direction of the normal out of each cell, as an angle counterclockwise from the x axis (ANGLDEGX), which MODFLOW 6 needs for horizontal anisotropy, the XT3D formulation, and specific discharge. triangle-utils calculates them for either grid (fig.~\ref{fig:connections}). In a Voronoi grid each face bisects the line between the two centers, so each distance is half the length of that line; the Voronoi grid of the valley has \voronoifaces\ faces between cells. In a triangle grid centered on centroids, the line between two centers is not perpendicular to their face, and the two distances, measured along the normal, sum to less than the distance between the centers. In a triangle grid centered on circumcenters, an obtuse triangle's circumcenter lies beyond its longest side, which makes the distance to that face negative, and a right triangle's lies on it, which makes the distance zero; triangle-utils reports such a face as an error, naming its two cells, rather than giving MODFLOW a distance that is not positive.

\begin{figure}[htbp]
  \centering
  \includegraphics{connections.pdf}
  \caption{Connections of one cell to the cells beside it in the example valley. \textit{A}, the triangles of the conforming Delaunay mesh, centered on their centroids. \textit{B}, the Voronoi grid of the same mesh (fig.~\ref{fig:voronoi}). Each shared face has a distance from the cell's center to the face, measured along the face's normal, a width, and the direction of its normal}
  \label{fig:connections}
\end{figure}

\section{Mapping Rivers, Boundaries, and Wells to Cells}

The Streamflow-Routing (SFR) package of MODFLOW 6 needs the cells a river passes through, in order, with the length of river in each, its reaches; the General-Head Boundary (GHB) package and other boundary packages need the cells along a boundary with the length of boundary in each, which sets the conductance of each cell's connection to the boundary. The Well (WEL) and Multi-Aquifer Well (MAW) packages need the cell of each well. triangle-utils finds all three in a triangle or Voronoi grid: the cells along a line given by its points, such as a river; the cells on the segments of a mesh with a marker, such as the valley's boundary, in the order the segments were digitized; and the cell containing each point, such as a well, or the lowest-numbered cell where a point is on a face or corner. A cell the line passes through twice appears twice along it; a part of a line along a face two cells share, as a river along a triangle's edge, is in the cell on the line's left.

A river network is often digitized as several lines, as a polyline shapefile gives it, with a line for each stream. Each line is added to the geometry with the same spacing and marker, and the cells along each are found in turn, so each stream's reaches are in order from its first point to its last. The valley's tributary enters across the east boundary and joins the river at the river's second bend from the north. The point the two lines share becomes one vertex of the mesh, and so the center of one cell, which holds a reach of each: the cell where the tributary's last reach flows into the river in the SFR package's connections. The cells on the segments with the streams' marker are the same cells, but they come back in three chains that break where the streams meet: the river above the confluence, the river below it, and the tributary. Stream reaches are found with the cells along each stream's line, and the cells on marked segments are used for boundaries, which do not branch.

Figure~\ref{fig:features} shows the cells of the streams and the boundary in the Voronoi grid of Figure~\ref{fig:voronoi}. The river's \riverlength~m pass through \rivercells\ cells and the tributary's \tributarylength~m through \tributarycells\ cells, each centered on the stream's vertices, and the boundary's \boundarylength~m through \boundarycells\ cells. A boundary cell's length is half the distance to the boundary vertex on each side of its own. Each well is the center of its ring, a vertex of the mesh, so its cell is centered on the well.

\begin{figure}[htbp]
  \centering
  \includegraphics{features.pdf}
  \caption{Cells of the streams, the boundary, and the wells in the Voronoi grid of the example valley (fig.~\ref{fig:voronoi}), with the windows of panels \textit{B}--\textit{D} (boxes). Each stream's cells alternate in color along it. \textit{A}, the valley. \textit{B}, at the confluence, where the tributary's last cell is one of the river's. \textit{C}, on the boundary at the valley's west corner. \textit{D}, around a well}
  \label{fig:features}
\end{figure}

\section{Adding Layers}

A MODFLOW 6 model is layered with the same horizontal grid in each layer. The vertex discretization (DISV) takes the top of each cell and the bottom of each cell in each layer, and an IDOMAIN value for each cell in each layer, which marks it active, outside the model (0), or a vertical pass-through cell ($-1$), through which the cells above and below it are connected as if it were not there. triangle-utils takes the same arrays and returns them for DISV, with the DISU arrays of the same layers. In DISU each active cell is a node, numbered layer by layer, with its top, bottom, and area. A node is connected to the active cells beside it in its layer, by the connections of Figure~\ref{fig:connections}, and to the active cells above and below it, through any pass-through cells between, by a vertical connection whose distance (CL12) is half the node's thickness and whose area (HWVA) is the cell's area, as MODFLOW 6 connects the layers of DISV. A cell outside the model breaks the vertical connection. MODFLOW 6 does not take pass-through cells in DISU, so they are left out of the nodes, and the DISV and DISU of the same layers give the same heads.

Elevations are often known at points, such as the land surface from a digital elevation model or the base of a unit from boreholes. triangle-utils interpolates values given at the mesh's vertices linearly to each cell's center, over the triangle that contains it; a Voronoi cell is centered on a vertex and takes that vertex's value. A cell's top is the bottom of the cell above it, as in MODFLOW. A layer whose bottom is not below the surface above it, as interpolated surfaces of a unit that pinches out can give, is an error that names the cell and layer, unless a minimum thickness is given; an active cell thinner than the minimum becomes a pass-through cell, so the flow between the units above and below it passes through. Cells outside the model are not checked.

The valley is given three layers (fig.~\ref{fig:layers}): an upper aquifer from the land surface, which rises from \landsurfacesouth~m at the south boundary by \landsurfacerise~m per kilometer to the north, to \upperbottom~m; a clay \claythickness~m thick at the west boundary that thins to nothing \clayedge~m east of it; and a lower aquifer to \lowerbottom~m. The land surface and the clay's base are given at the mesh's vertices. With a minimum thickness of \minimumthickness~m, \passthroughcells\ cells of the clay pass flow between the aquifers. The DISV grid has \disvcells\ cells in its three layers; the DISU grid has \disunodes\ nodes, one for each active cell, and \disuconnections\ entries in its connection arrays, a node's connection to itself included.

\begin{figure}[htbp]
  \centering
  \includegraphics{layers.pdf}
  \caption{Layers of the Voronoi grid of the example valley (fig.~\ref{fig:voronoi}). \textit{A}, the thickness of the clay, the second layer, with the cells thinner than \minimumthickness~m, which become pass-through cells, and the line of section A--A$'$ at y = \sectiony~m. \textit{B}, the three layers along section A--A$'$; each column is a cell the section crosses, as wide as the length of section in it, and the pass-through cells of the clay are too thin to show}
  \label{fig:layers}
\end{figure}

\section{Using triangle-utils}

The \texttt{triangle-utils} program reads and writes the files of Triangle. The \texttt{generate} command meshes the segments, vertices, holes, and regions of a .poly file, such as one describing the boundary, river, and lake of a valley, the \texttt{refine} command refines the mesh in a .node file and an .ele file, with the segments of a .poly file and the area of each triangle in an .area file if they exist, the \texttt{improve} command moves the vertices of such a mesh, the \texttt{disv} command writes its Voronoi grid, or its triangles, as a MODFLOW 6 DISV file, and the \texttt{cells} command prints the cells of the grid that points contain or that a line or the segments with a marker pass through. A vertex with marker 1 in a .node file, as \texttt{generate} writes for the points of the .poly file, such as wells, stays in place:

\begin{Verbatim}[fontsize=\small,commandchars=\\\{\}]
triangle-utils generate -a \maximumareavalue -q \minimumangle -s \boundaryspacingvalue -D valley
triangle-utils refine -A \maximumangle -D valley.1
triangle-utils improve -A \improveangle valley.2
triangle-utils check valley.3
triangle-utils disv valley.2
triangle-utils cells -l river.txt valley.2
triangle-utils cells -l tributary.txt valley.2
triangle-utils cells -m \boundarymarker valley.2
triangle-utils cells -p wells.txt valley.2
\end{Verbatim}

\noindent Each command writes its mesh to files named as Triangle names them, valley.1 for valley and valley.2 for valley.1, and \texttt{disv} writes the Voronoi grid of the refined mesh to valley.2.disv, as Figure~\ref{fig:voronoi} shows; \texttt{triangle-utils generate -h}, \texttt{refine -h}, \texttt{improve -h}, \texttt{disv -h}, and \texttt{cells -h} list the options of each. The files river.txt, tributary.txt, and wells.txt give the x and y of each of a stream's points and of each well, one point to a line, such as the coordinates of RIVER, TRIBUTARY, and WELLS in valley.py below; each stream is mapped by a command of its own.

From Python, the triangle-utils module in the python folder of the repository calls the library directly, and meshes and grids are returned as NumPy arrays. The script valley.py in the doc/UserGuide folder builds the example of this guide in steps, printed here in turn; the script make\_figures.py beside it uses its meshes and grid to make the figures. First come the coordinates of the boundary, river, tributary, lake, and wells, with the streams as a list of lines:

\VerbatimInput[fontsize=\small,firstline=\valleycoordinatesfirst,lastline=\valleycoordinateslast]{valley.py}

\noindent Then come the spacings of the points along each feature, the largest triangle areas, the angles, the elevations of the layers, and the attribute of each region and marker of each feature's segments:

\VerbatimInput[fontsize=\small,firstline=\valleyvaluesfirst,lastline=\valleyvalueslast]{valley.py}

\noindent The geometry is built from them: each feature with the spacing of its points and the marker of its segments, a point inside each region with its attribute and area limit, and a ring around each well:

\VerbatimInput[fontsize=\small,firstline=\valleygeometryfirst,lastline=\valleygeometrylast]{valley.py}

\noindent The geometry is meshed (fig.~\ref{fig:generated}), the mesh is refined by angle (fig.~\ref{fig:refined}), and the refined mesh is improved (fig.~\ref{fig:improved}):

\VerbatimInput[fontsize=\small,firstline=\valleymeshingfirst,lastline=\valleymeshinglast]{valley.py}

\noindent The Voronoi grid (fig.~\ref{fig:voronoi}) is built from the geometry meshed and refined again with conforming Delaunay triangles, and its vertices and cells are returned in the form FloPy's DISV package takes:

\VerbatimInput[fontsize=\small,firstline=\valleygridfirst,lastline=\valleygridlast]{valley.py}

\noindent The connections between the grid's cells (fig.~\ref{fig:connections}) are calculated from the grid:

\VerbatimInput[fontsize=\small,firstline=\valleyconnectionsfirst,lastline=\valleyconnectionslast]{valley.py}

\noindent Each stream's, the boundary's, and the wells' cells (fig.~\ref{fig:features}) are found in the grid, each stream's from its own line and the boundary's from the segments of the mesh with the boundary's marker:

\VerbatimInput[fontsize=\small,firstline=\valleyfeaturesfirst,lastline=\valleyfeatureslast]{valley.py}

\noindent The grid is given three layers (fig.~\ref{fig:layers}), from elevations at the mesh's vertices, and the DISV and DISU arrays are returned as the arguments of FloPy's packages, used as \texttt{flopy.mf6.ModflowGwfdisv(model, **disv)}:

\VerbatimInput[fontsize=\small,firstline=\valleylayersfirst,lastline=\valleylayerslast]{valley.py}

Features drawn in a GIS are read with \texttt{read\_features}, from a shapefile, GeoPackage, or GeoJSON file, a GeoDataFrame, or shapely geometries; it needs geopandas, installed with the module as triangle-utils[gis]. It returns each part of each feature, with the points the functions above take, the row of the table it came from, and its marker, spacing, region attribute, and area limit, each given as a number or as the name of the column that holds it. A polygon's holes are returned as hole polygons, and every polygon's points run counterclockwise, as in valley.py, although a shapefile stores an outer boundary clockwise. The mesh is in plain x and y, so the features must be in a projected coordinate system; a geographic coordinate system is an error, and features without one are taken to be projected. With the valley's boundary, streams, and lake in shapefiles whose marker and spacing columns hold the values above, a point inside each region in a shapefile with zone and area columns, and the wells in a shapefile of points, the geometry is built by

\begin{Verbatim}[fontsize=\small]
columns = {"marker": "marker", "spacing": "spacing"}
geometry = tu.Geometry()
for path in ["boundary.shp", "streams.shp", "lake.shp"]:
    geometry.add_features(tu.read_features(path, **columns))
regions = tu.read_features("regions.shp", attribute="zone", area_limit="area")
geometry.add_features(regions, points="regions")
for well in tu.read_features("wells.shp"):
    x, y = well.points[0]
    geometry.add_ring(
        x,
        y,
        radius=WELL_RADIUS,
        point_count=WELL_POINTS,
        attribute=WELL_REGION,
        area_limit=WELL_AREA,
        marker=WELL_MARKER,
    )
\end{Verbatim}

\noindent and meshes to the same mesh as the coordinates in valley.py. Each stream read from streams.shp has the points \texttt{cells\_along\_line} takes.

\noindent The same functions are available to C, C++, and Fortran programs through the C interface declared in triangle\_utils.h.

\REFSECTION
\bibliographystyle{usgs}
\bibliography{triangle-utils}

\end{document}
